% Gemeinsamer Stil aller lenscape-Abbildungen (TikZ/pgfplots).
% Einbinden mit \input{../lenscape-style} aus einem Abbildungsordner.

% Schrift: Computer Modern (in jeder TeX-Installation vorhanden, passt zu LaTeX-Dokumenten).
% Für Latin Modern: \usepackage[T1]{fontenc}\usepackage{lmodern}
\usepackage{amsmath}
\usepackage{pgfplots}
\pgfplotsset{compat=1.18}
\usepgfplotslibrary{groupplots}
\usetikzlibrary{calc}

% ---------------------------------------------------------------- Farben
\definecolor{lsBlue}{HTML}{2A6FDB}
\definecolor{lsOrange}{HTML}{E8833A}
\definecolor{lsGrey}{HTML}{6B7280}
\definecolor{lsGrid}{HTML}{E5E7EB}
\definecolor{lsInk}{HTML}{1F2937}
\definecolor{lsRed}{HTML}{C2413A}
\definecolor{lsGreen}{HTML}{2E8B57}
% sechs gut unterscheidbare Farben für Einzelobjekte
\definecolor{lsC1}{HTML}{2A6FDB}
\definecolor{lsC2}{HTML}{E8833A}
\definecolor{lsC3}{HTML}{2E8B57}
\definecolor{lsC4}{HTML}{9B59B6}
\definecolor{lsC5}{HTML}{C2413A}
\definecolor{lsC6}{HTML}{8C6D31}

% Deutsches Dezimalkomma auch für Zahlen außerhalb von Achsen
\pgfkeys{/pgf/number format/use comma}

% deutsche Anführungszeichen ohne babel/T1
\providecommand{\glqq}{\raisebox{-0.9ex}{''}}
\providecommand{\grqq}{``}

% ---------------------------------------------------------------- Kürzel
\newcommand{\Msun}{\ensuremath{M_\odot}}
\newcommand{\thetaE}{\ensuremath{\theta_\mathrm{E}}}

% ---------------------------------------------------------------- Achsenstil
\pgfplotsset{
  lenscape/.style={
    width=8.6cm, height=5.8cm,
    axis line style={lsGrey!80},
    tick style={lsGrey!80},
    every tick label/.append style={font=\small, text=lsInk},
    label style={font=\small, text=lsInk},
    title style={font=\small\bfseries, text=lsInk},
    grid=major, major grid style={lsGrid, line width=0.4pt},
    legend style={font=\footnotesize, draw=none, fill=white, fill opacity=0.85, text opacity=1},
    legend cell align=left,
    /pgf/number format/use comma,
    scaled ticks=false,
  },
  lenscape wide/.style={lenscape, width=8.2cm, height=6.2cm},
  % Datenpunkte
  dataBlue/.style={only marks, mark=*, mark size=1.4pt, color=lsBlue, fill opacity=0.85,
                   error bars/error bar style={lsBlue!60, line width=0.4pt}},
  dataOrange/.style={only marks, mark=square*, mark size=1.3pt, color=lsOrange, fill opacity=0.85,
                     error bars/error bar style={lsOrange!60, line width=0.4pt}},
  % Konfidenzbereiche (68 % / 95 %) für echte Daten und Prognose
  slacs68/.style={fill=lsBlue, fill opacity=0.55, draw=lsBlue!70!black, line width=0.5pt},
  slacs95/.style={fill=lsBlue, fill opacity=0.20, draw=lsBlue!70!black, line width=0.5pt},
  forecast68/.style={fill=lsOrange, fill opacity=0.60, draw=lsOrange!70!black, line width=0.5pt},
  forecast95/.style={fill=lsOrange, fill opacity=0.22, draw=lsOrange!70!black, line width=0.5pt},
  lens68/.style={fill=lsBlue, fill opacity=0.55, draw=lsBlue!70!black, line width=0.5pt},
  lens95/.style={fill=lsBlue, fill opacity=0.20, draw=lsBlue!70!black, line width=0.5pt},
  reference/.style={lsInk, line width=0.6pt},
}

% Vergleichsbänder (frühes vs. spätes Universum), für \fill innerhalb von Achsen
\tikzset{
  bandEarly/.style={fill=lsRed, fill opacity=0.15, draw=none},
  bandLate/.style={fill=lsGreen, fill opacity=0.15, draw=none},
}
